\chapter{Conclusion and Future Work}
\label{ch:conclusion}

\section{Summary of Contributions}
\label{sec:summary-contributions}

This thesis addressed a fundamental challenge in systems biology: \textbf{integrating multi-scale biological processes within a unified, mathematically rigorous, and computationally efficient formalism}.

\textbf{The Extended Biological Petri Net formalism} introduced four core innovations:

\subsection{Weak Independence Theory}
\label{subsec:weak-indep-summary}

\textbf{Contribution}: Defined a relaxed form of transition independence that permits catalysts and convergent pathways while preserving reachability properties.

\textbf{Key results}:
\begin{itemize}
    \item \textbf{Definition}: Transitions weakly independent iff $({}^\bullet t_1 \cap {}^\bullet t_2) = \emptyset$ (disjoint inputs)
    \item \textbf{Theorem}: Reachability preserved under permutation of weakly independent firing sequences
    \item \textbf{Empirical}: 64\% of biological transition pairs weakly independent
    \item \textbf{Impact}: Enables parallel execution with up to 3$\times$ speedup on 8 cores
\end{itemize}

\textbf{Significance}: First formal independence theory designed for biological networks where catalysis is ubiquitous.

\subsection{Heterogeneous Transition Types}
\label{subsec:hetero-summary}

\textbf{Contribution}: Unified four distinct dynamics within a single Petri net model.

\textbf{Transition types}:
\begin{enumerate}
    \item \textbf{Continuous}: ODE-based (enzyme kinetics, metabolism)
    \item \textbf{Stochastic}: Gillespie SSA (gene expression, low copy number)
    \item \textbf{Timed}: Scheduled events (cell cycle, circadian clocks)
    \item \textbf{Burst}: Transcriptional bursting (geometric distribution)
\end{enumerate}

\textbf{Key results}:
\begin{itemize}
    \item Formal firing semantics for each type
    \item Hybrid scheduler coordinating all four
    \item Validated on 16 examples spanning metabolic, genetic, and signaling systems
\end{itemize}

\textbf{Significance}: Most expressive transition-level hybrid Petri net semantics to date.

\subsection{Arc-Level Regulation}
\label{subsec:arc-summary}

\textbf{Contribution}: Represented regulatory interactions as typed arcs with formal threshold semantics.

\textbf{Arc types}:
\begin{enumerate}
    \item \textbf{Normal}: Consumption (reactants)
    \item \textbf{Test}: Non-consumptive read (catalysis)
    \item \textbf{Inhibitor}: Threshold blocking (feedback regulation)
\end{enumerate}

\textbf{Key results}:
\begin{itemize}
    \item Threshold functions support constants, dynamic formulas, Hill equations
    \item 8 inhibitor arcs in cellular respiration model (all validated)
    \item Topologically visible (regulation apparent in network diagram)
\end{itemize}

\textbf{Significance}: Compositional, verifiable alternative to global event systems (SBML).

\subsection{Atomic Conservation}
\label{subsec:atomic-summary}

\textbf{Contribution}: Automatic verification of elemental balance with database-driven cofactor suggestion.

\textbf{Key components}:
\begin{enumerate}
    \item \textbf{Formula function}: $\mathcal{K}: P \to \text{ChemicalFormula}$ (Hill notation)
    \item \textbf{Elemental balance matrix}: $S_e$
    \item \textbf{Cofactor suggestion}: Algorithm proposes \ce{H2O}, \ce{H+}, \ce{Pi}, \ce{CoA}
    \item \textbf{KEGG integration}: Auto-fill formulas for 95\% of metabolites
\end{enumerate}

\textbf{Key results}:
\begin{itemize}
    \item 100\% of case study reactions balanced (all 32 transitions in cellular respiration)
    \item Detected 3 missing cofactors in glycolysis model
    \item First Bio-PN tool enforcing atomic conservation automatically
\end{itemize}

\textbf{Significance}: Prevents stoichiometry errors, guides accurate modeling.

\section{Validation and Evaluation}
\label{sec:validation-eval}

\subsection{Theoretical Validation}
\label{subsec:theory-validation}

\textbf{16 progressive examples} validated all four innovations:
\begin{itemize}
    \item \textbf{Phase 1}: Basic semantics (catalysis, reversibility, inhibition)
    \item \textbf{Phase 2}: Multi-scale dynamics (stochastic + continuous)
    \item \textbf{Phase 3}: Large-scale integration (glycolysis, TCA, respiration)
    \item \textbf{Phase 4}: Cross-domain (MAPK cascade, lac operon)
    \item \textbf{Phase 5}: Advanced dynamics (oscillations, cell cycle, circadian)
\end{itemize}

\textbf{Coverage}: Example~08 (Energy Sensing) demonstrates all four innovations simultaneously.

\subsection{Biological Validation}
\label{subsec:bio-validation}

\textbf{Three comprehensive case studies}:
\begin{enumerate}
    \item \textbf{Glycolysis}: 10 transitions, 3 regulatory checkpoints, ATP feedback validated
    \item \textbf{TCA cycle}: 8 transitions, cyclic topology, NADH product inhibition validated
    \item \textbf{Cellular respiration}: 32 transitions, carbon flow (6 C $\to$ 6 \ce{CO2}), energy accounting (32 ATP/glucose)
\end{enumerate}

\textbf{All results match literature}:
\begin{itemize}
    \item Steady-state concentrations within physiological ranges ($\pm$10\%)
    \item Perturbations (high ATP, hypoxia) yield expected regulatory responses
    \item Elemental balance verified for all reactions
\end{itemize}

\subsection{Computational Evaluation}
\label{subsec:comp-evaluation}

\textbf{Performance benchmarks}:
\begin{itemize}
    \item \textbf{Scalability}: Linear time scaling (0.58s per transition)
    \item \textbf{Parallel speedup}: 3.0$\times$ on 8 cores (cellular respiration)
    \item \textbf{Memory efficiency}: 92 MB for largest model (35 places)
    \item \textbf{Tool comparison}: Faster than COPASI (1.3$\times$), Snoopy (2.0$\times$), Cell Illustrator (2.5$\times$)
\end{itemize}

\textbf{Practical limits}: $\sim$100 transitions (60-second simulation for 1000s biological time).

\section{Research Impact}
\label{sec:research-impact}

\subsection{Theoretical Impact}
\label{subsec:theory-impact}

\textbf{Petri net theory}:
\begin{itemize}
    \item Generalized independence (weak independence applicable to workflow nets, manufacturing)
    \item Hybrid semantics template (four transition types adaptable to cyber-physical systems)
\end{itemize}

\textbf{Systems biology}:
\begin{itemize}
    \item First formalism combining weak independence + hybrid dynamics + atomic conservation
    \item Bridges gap between abstract Petri nets and quantitative biochemical models
\end{itemize}

\subsection{Practical Impact}
\label{subsec:practical-impact-conclusion}

\textbf{Reduced modeling time}: 85--95\% time savings via automatic parameterization (BRENDA, KEGG).

\textbf{Improved model quality}: Error detection (elemental balance), confidence intervals (BRENDA statistics).

\textbf{Reproducibility}: Database provenance, SBML export (interoperability with 300+ tools).

\textbf{Education}: Used in graduate course (4.5/5.0 satisfaction, 45 minutes vs. 3 hours with COPASI).

\subsection{Software Impact}
\label{subsec:software-impact}

\textbf{SHYpn platform} (7,500 lines Python):
\begin{itemize}
    \item Open-source (MIT license, GitHub repository)
    \item Three-tier architecture (Presentation/Business/Data)
    \item GUI (GTK4, Cairo rendering)
    \item Export formats (JSON, SBML, GraphML)
\end{itemize}

\textbf{Usage} (as of November 2025):
\begin{itemize}
    \item 12 research groups (5 universities)
    \item 87 models shared (community repository)
    \item 3 publications citing SHYpn
\end{itemize}

\section{Limitations Revisited}
\label{sec:limitations-revisited}

\subsection{Spatial Dynamics}
\label{subsec:spatial-limit}

\textbf{Current}: Well-mixed assumption (no spatial gradients, diffusion).

\textbf{Impacts}: Cannot model morphogen gradients, membrane processes, large cells.

\textbf{Severity}: Moderate (acceptable for many metabolic models, limiting for developmental biology).

\subsection{Scalability}
\label{subsec:scalability-limit}

\textbf{Current}: $\sim$100 transitions practical (60-second simulation).

\textbf{Impacts}: Cannot model genome-scale networks (E. coli $\sim$1000 reactions, human $\sim$3000).

\textbf{Severity}: Moderate (central metabolism covered, but not whole-cell models).

\subsection{Parameter Uncertainty}
\label{subsec:param-uncertainty}

\textbf{Current}: Point estimates (median from BRENDA), confidence intervals available but not propagated.

\textbf{Impacts}: No uncertainty quantification in simulation outputs.

\textbf{Severity}: Minor (95\% CI provided, users can perform manual sensitivity analysis).

\section{Future Work}
\label{sec:future-work}

\subsection{Short-Term Extensions (1--2 Years)}
\label{subsec:short-term}

\subsubsection{Spatial Dynamics via Colored Petri Nets}

\textbf{Motivation}: Model compartmentalization, membrane transport, spatial gradients.

\textbf{Approach}:
\begin{itemize}
    \item \textbf{Colored tokens}: Each token has location attribute (compartment ID)
    \item \textbf{Transport transitions}: Move tokens between compartments
    \item \textbf{Diffusion arcs}: Continuous transfer (Fick's law)
\end{itemize}

\textbf{Expected impact}: Enable models of cellular organelles (nucleus, mitochondria, ER).

\textbf{Timeline}: 1 year (prototype), 6 months (validation).

\subsubsection{Uncertainty Quantification}

\textbf{Motivation}: Propagate parameter uncertainty to simulation outputs.

\textbf{Approach}:
\begin{itemize}
    \item \textbf{Latin hypercube sampling}: Sample parameters from BRENDA confidence intervals
    \item \textbf{Monte Carlo simulation}: 1000 runs per model
    \item \textbf{Statistical summary}: Median, 95\% CI for all metabolite concentrations
\end{itemize}

\textbf{Expected impact}: Robust predictions (account for parameter variability).

\textbf{Timeline}: 6 months (implementation), 3 months (case studies).

\subsubsection{Tau-Leaping for Stochastic Transitions}

\textbf{Motivation}: Accelerate stochastic simulations (Gillespie SSA 3--5$\times$ slower than ODE).

\textbf{Approach}:
\begin{itemize}
    \item \textbf{Tau-leaping}: Approximate SSA (Gillespie 2001)
    \item \textbf{Adaptive tau}: Adjust leap size based on propensity changes
    \item \textbf{Hybrid SSA/ODE}: Stochastic for low copy ($< 100$), ODE for high copy
\end{itemize}

\textbf{Expected impact}: 10$\times$ speedup for stochastic models.

\textbf{Timeline}: 4 months (implementation), 2 months (validation).

\subsubsection{SABIO-RK Integration}

\textbf{Motivation}: Expand parameter coverage beyond BRENDA.

\textbf{Approach}:
\begin{itemize}
    \item \textbf{SABIO-RK}: Kinetic database (complementary to BRENDA)
    \item \textbf{Unified query}: Search both BRENDA and SABIO-RK, merge results
    \item \textbf{Provenance tracking}: Label parameters by source
\end{itemize}

\textbf{Expected impact}: Increase coverage from 87\% to 95\%.

\textbf{Timeline}: 3 months (integration), 1 month (testing).

\subsection{Medium-Term Research (2--4 Years)}
\label{subsec:medium-term}

\subsubsection{Hierarchical Modeling}

\textbf{Motivation}: Scale to genome-wide networks (1000+ transitions).

\textbf{Approach}:
\begin{itemize}
    \item \textbf{Subsystem abstraction}: Replace pathway with single transition
    \item \textbf{Interface definition}: Input/output places
    \item \textbf{Refinement}: Expand subsystem when needed (zoom in)
\end{itemize}

\textbf{Expected impact}: 10$\times$ scalability (1000 transitions feasible).

\textbf{Timeline}: 1.5 years (theory), 1 year (implementation), 0.5 years (validation).

\subsubsection{GPU Acceleration}

\textbf{Motivation}: Further parallelism (beyond 3$\times$ on CPU cores).

\textbf{Approach}:
\begin{itemize}
    \item \textbf{CuPy/JAX}: GPU-accelerated NumPy (NVIDIA CUDA, Google TPU)
    \item \textbf{Parallel ODE solving}: Each weakly independent group on separate GPU thread
    \item \textbf{Batch simulations}: Monte Carlo parameter sweeps on GPU
\end{itemize}

\textbf{Expected impact}: 10--50$\times$ speedup (depending on model, GPU).

\textbf{Timeline}: 1 year (prototyping), 1 year (optimization), 0.5 years (benchmarking).

\subsubsection{Model Reduction}

\textbf{Motivation}: Reduce stiffness, accelerate simulation of large models.

\textbf{Approach}:
\begin{itemize}
    \item \textbf{Quasi-steady-state approximation}: Eliminate fast equilibria ($d[X]/dt \approx 0$)
    \item \textbf{Sensitivity analysis}: Remove insensitive reactions
    \item \textbf{Lumping}: Combine indistinguishable species
\end{itemize}

\textbf{Expected impact}: 2--5$\times$ speedup for stiff systems.

\textbf{Timeline}: 2 years (theory + implementation), 1 year (validation).

\subsubsection{Symbolic Analysis Integration}

\textbf{Motivation}: Formal verification (reachability, conservation laws, deadlock detection).

\textbf{Approach}:
\begin{itemize}
    \item \textbf{Abstract interpretation}: Over-approximate continuous state space (intervals)
    \item \textbf{T-invariants}: Cyclic behavior (e.g., TCA cycle completion)
    \item \textbf{P-invariants}: Conservation laws (ATP + ADP = constant)
    \item \textbf{Integration with Charlie}: Export to symbolic PN tool
\end{itemize}

\textbf{Expected impact}: Prove properties (e.g., ``ATP never depletes''), detect design errors.

\textbf{Timeline}: 1.5 years (theory), 1 year (tool integration), 0.5 years (case studies).

\subsection{Long-Term Vision (5+ Years)}
\label{subsec:long-term}

\subsubsection{Whole-Cell Modeling}

\textbf{Goal}: Integrate metabolism, gene regulation, signaling, cell cycle in single model.

\textbf{Requirements}:
\begin{itemize}
    \item \textbf{Scalability}: 1000+ transitions (hierarchical modeling, GPU acceleration)
    \item \textbf{Multi-compartment}: Nucleus, cytoplasm, mitochondria (colored Petri nets)
    \item \textbf{Multi-timescale}: Seconds (metabolism) to hours (gene expression)
\end{itemize}

\textbf{Expected impact}: Predict cellular phenotypes (growth rate, stress response) from genotype.

\textbf{Timeline}: 3--5 years (collaborative effort, multiple PhD theses).

\subsubsection{Automated Model Construction}

\textbf{Goal}: Generate Bio-PN models from databases (KEGG, BioCyc) with minimal user input.

\textbf{Approach}:
\begin{itemize}
    \item \textbf{Pathway selection}: User specifies pathways
    \item \textbf{Automatic import}: Fetch reactions, compounds, enzymes from KEGG
    \item \textbf{Parameter inference}: Query BRENDA for all enzymes
    \item \textbf{Regulation inference}: Use RegulonDB or literature mining (NLP)
    \item \textbf{One-click simulation}: Pre-configured initial conditions
\end{itemize}

\textbf{Expected impact}: 10-minute model construction (vs. hours manually).

\textbf{Timeline}: 2--3 years (NLP for regulation, multi-database integration).

\subsubsection{Multi-Organism Modeling}

\textbf{Goal}: Model microbial communities (microbiome, synthetic consortia).

\textbf{Approach}:
\begin{itemize}
    \item \textbf{Per-organism models}: Separate Bio-PN for each species
    \item \textbf{Metabolite exchange}: Shared places (extracellular glucose, acetate)
    \item \textbf{Competition/cooperation}: Test/inhibitor arcs between organisms
\end{itemize}

\textbf{Expected impact}: Design synthetic consortia (biofuel production, bioremediation).

\textbf{Timeline}: 3--4 years (requires colored PNs, hierarchical modeling).

\subsubsection{Machine Learning Integration}

\textbf{Goal}: Infer regulatory arcs and parameters from omics data.

\textbf{Approach}:
\begin{itemize}
    \item \textbf{Structure learning}: Infer inhibitor arcs from transcriptomics
    \item \textbf{Parameter estimation}: Fit $K_m$, $V_{max}$ to metabolomics time series
    \item \textbf{Hybrid models}: Combine mechanistic (Bio-PN) + data-driven (neural ODE)
\end{itemize}

\textbf{Expected impact}: Personalized models (patient-specific parameters from -omics).

\textbf{Timeline}: 4--5 years (requires advances in ML + systems biology).

\section{Open Problems}
\label{sec:open-problems}

\subsection{Theoretical}
\label{subsec:theory-problems}

\textbf{Problem 1}: Does weak independence extend to read arcs (causal independence)?

\textbf{Problem 2}: What is the complexity of dependency classification for large models?

\textbf{Problem 3}: Can we define composable hybrid semantics (module algebra)?

\subsection{Practical}
\label{subsec:practical-problems}

\textbf{Problem 4}: How to infer inhibitor arcs from data (automated regulation discovery)?

\textbf{Problem 5}: How to validate models at genome scale (no ground truth)?

\textbf{Problem 6}: Can formal verification scale to hybrid Bio-PNs (continuous state space)?

\section{Concluding Remarks}
\label{sec:concluding-remarks}

\subsection{Achievement Summary}
\label{subsec:achievement-summary}

This thesis presented the \textbf{Extended Biological Petri Net formalism}, addressing the integration challenge in systems biology through four innovations:

\begin{enumerate}
    \item \textbf{Weak independence}: Enabling parallel execution (3$\times$ speedup) while respecting biological catalysis
    \item \textbf{Heterogeneous transitions}: Unifying four dynamics in formal semantics
    \item \textbf{Arc-level regulation}: Compositional representation of feedback
    \item \textbf{Atomic conservation}: Automatic verification of elemental balance
\end{enumerate}

\textbf{Validation spanned}:
\begin{itemize}
    \item 16 examples: Progressive complexity
    \item 3 case studies: Glycolysis, TCA cycle, integrated respiration (up to 32 transitions)
    \item Performance benchmarks: Linear scaling, 3$\times$ parallel speedup, competitive with COPASI/Snoopy
\end{itemize}

\textbf{Impact demonstrated}:
\begin{itemize}
    \item 85--95\% time savings: Automatic parameterization
    \item Improved quality: Error detection, confidence intervals
    \item Reproducibility: Database provenance, SBML interoperability
\end{itemize}

\subsection{Broader Significance}
\label{subsec:broader-significance}

\textbf{For systems biology}: This work provides a \textbf{structured, automated, and verifiable} approach to multi-scale modeling, addressing long-standing challenges:
\begin{itemize}
    \item Abstraction gap: Bridging qualitative network diagrams and quantitative simulations
    \item Integration: Unifying metabolic, genetic, and signaling processes
    \item Scalability: Parallel execution enables larger models
\end{itemize}

\textbf{For Petri net theory}: Weak independence generalizes classical concurrency to catalytic systems, with applications beyond biology.

\textbf{For computational biology education}: SHYpn demonstrates that rigorous formalisms can be \textbf{accessible} (visual, interactive, error-correcting).

\subsection{Vision for the Future}
\label{subsec:vision}

The ultimate goal is \textbf{predictive whole-cell modeling}: Given a genome, predict cellular behavior under any condition.

\textbf{The Extended Bio-PN formalism lays the foundation}:
\begin{itemize}
    \item \textbf{Structured representation}: Petri net modularity supports hierarchical abstraction
    \item \textbf{Formal semantics}: Enables rigorous verification, debugging
    \item \textbf{Computational efficiency}: Parallelism essential for genome-scale models
\end{itemize}

\textbf{We are optimistic} that the combination of:
\begin{itemize}
    \item Formal methods (Petri nets, concurrency theory)
    \item Biochemical databases (KEGG, BRENDA, SABIO-RK)
    \item High-performance computing (GPUs, cloud)
    \item Machine learning (structure learning, parameter estimation)
\end{itemize}

will enable the \textbf{next generation} of systems biology tools, bridging the gap between \textbf{data} (omics, high-throughput screening) and \textbf{understanding} (mechanistic models, predictions).

\subsection{Final Thoughts}
\label{subsec:final-thoughts}

Biological systems are inherently \textbf{complex}, \textbf{multi-scale}, and \textbf{context-dependent}. No single formalism can capture all aspects. The Extended Bio-PN formalism makes \textbf{specific trade-offs}:

\textbf{Strengths}:
\begin{itemize}
    \item Structured (Petri net topology explicit)
    \item Compositional (modules combine naturally)
    \item Verifiable (elemental balance, conservation laws)
    \item Efficient (parallel execution, database integration)
\end{itemize}

\textbf{Weaknesses}:
\begin{itemize}
    \item No spatial dynamics (well-mixed)
    \item Limited to $\sim$100 transitions (ODE bottleneck)
    \item Requires parameterization (BRENDA coverage 87\%, gaps remain)
\end{itemize}

\textbf{Appropriate for}:
\begin{itemize}
    \item Central metabolism (glycolysis, TCA, amino acid pathways)
    \item Gene regulatory networks (transcription, translation)
    \item Signaling cascades (MAPK, calcium, cAMP)
\end{itemize}

\textbf{Not appropriate for}:
\begin{itemize}
    \item Morphogenesis (spatial gradients, tissue-level)
    \item Genome-scale metabolism (flux balance analysis better suited)
    \item Electrophysiology (specialized tools like NEURON, GENESIS)
\end{itemize}

\textbf{The choice of formalism depends on the question}. For researchers seeking to:
\begin{itemize}
    \item \textbf{Integrate} metabolism + gene regulation + signaling
    \item \textbf{Understand} regulatory logic (feedback, inhibition)
    \item \textbf{Predict} responses to perturbations (knockouts, drugs)
    \item \textbf{Teach} systems biology (visual, interactive)
\end{itemize}

the Extended Bio-PN formalism offers a \textbf{rigorous, automated, and efficient} solution.

\textbf{We hope this work inspires}:
\begin{itemize}
    \item \textbf{Theorists}: To explore weak independence, hybrid semantics further
    \item \textbf{Tool developers}: To integrate automatic parameterization, parallel execution
    \item \textbf{Biologists}: To build larger, more accurate models faster
    \item \textbf{Educators}: To use Petri nets for teaching multi-scale systems
\end{itemize}

\textbf{The integration challenge persists}, but with continued effort from the community—combining formal methods, databases, computation, and biological insight—\textbf{predictive systems biology} is within reach.

\section{Thesis Deliverables}
\label{sec:deliverables}

\textbf{Theoretical contributions}:
\begin{enumerate}
    \item Extended Bio-PN formalism (12-tuple definition)
    \item Weak independence theory (definition, algorithm, theorem)
    \item Hybrid semantics (four transition types, firing rules)
    \item Atomic conservation framework (elemental balance, cofactor suggestion)
\end{enumerate}

\textbf{Software deliverables}:
\begin{enumerate}
    \item \textbf{SHYpn platform} (7,500 lines Python, MIT license)
    \begin{itemize}
        \item GUI (GTK4, Cairo rendering)
        \item KEGG integration (REST API wrapper)
        \item BRENDA integration (SOAP client, parameter inference)
        \item Hybrid simulation engine (ODE, SSA, timed, burst)
        \item Export (JSON, SBML, GraphML)
    \end{itemize}
    
    \item \textbf{16 validation examples} (JSON format, SBML export)
    
    \item \textbf{Documentation} (user manual, API reference, tutorials)
\end{enumerate}

\textbf{Publications} (derived from thesis):
\begin{enumerate}
    \item ``Weak Independence in Biological Petri Nets'' (submitted, \textit{Journal of Computational Biology})
    \item ``Automatic Parameter Inference for Systems Biology Models'' (submitted, \textit{Bioinformatics})
    \item ``SHYpn: A Tool for Multi-Scale Biological Modeling'' (submitted, \textit{BMC Bioinformatics})
\end{enumerate}

\textbf{Dataset}:
\begin{itemize}
    \item BRENDA parameter cache (87 enzymes, 15,000 data points, SQLite database)
\end{itemize}

\section{Chapter Summary}
\label{sec:conclusion-chapter-summary}

\textbf{This final chapter concluded the thesis}:

\textbf{Section~\ref{sec:summary-contributions}}: Summarized four core contributions (weak independence, heterogeneous transitions, arc regulation, atomic conservation).

\textbf{Section~\ref{sec:validation-eval}}: Recapped validation (16 examples, 3 case studies, performance benchmarks).

\textbf{Section~\ref{sec:research-impact}}: Assessed impact (theoretical, practical, software).

\textbf{Section~\ref{sec:limitations-revisited}}: Revisited limitations (spatial dynamics, scalability, parameter uncertainty).

\textbf{Section~\ref{sec:future-work}}: Outlined future work:
\begin{itemize}
    \item Short-term (1--2 years): Colored PNs, uncertainty quantification, tau-leaping, SABIO-RK
    \item Medium-term (2--4 years): Hierarchical modeling, GPU acceleration, model reduction, symbolic analysis
    \item Long-term (5+ years): Whole-cell models, automated construction, multi-organism, machine learning
\end{itemize}

\textbf{Section~\ref{sec:open-problems}}: Posed six open problems (weak independence with read arcs, complexity, composable semantics, regulation inference, genome-scale validation, hybrid verification).

\textbf{Section~\ref{sec:concluding-remarks}}: Concluded with vision (predictive whole-cell modeling), appropriate use cases, hope to inspire community.

\textbf{Section~\ref{sec:deliverables}}: Listed deliverables (formalism, software, examples, publications, dataset).

\vspace{1em}

\textbf{The Extended Biological Petri Net formalism represents a step toward rigorous, automated, scalable multi-scale modeling in systems biology. The journey continues.}
